% Luc Destombe --- licence CC BY-NC-SA 4.0
% https://creativecommons.org/licenses/by-nc-sa/4.0/deed.fr
% Fichier autonome : compiler avec pdflatex, deux fois (signets, nombre de pages).
\def\lyceeversion{1.0}
\documentclass[11pt,a4paper]{article}
\makeatletter
\newif\ifeleve
\elevefalse

\newif\iflycee@cm
\newif\iflycee@palatino

\RequirePackage[utf8]{inputenc}
\RequirePackage[T1]{fontenc}
\RequirePackage[french]{babel}
\RequirePackage{etoolbox}

\newcommand{\ShorthandsOff}{\@ifundefined{shorthandoff}{}{\shorthandoff{;:!?}}}
\newcommand{\ShorthandsOn}{\@ifundefined{shorthandon}{}{\shorthandon{;:!?}}}
\ShorthandsOff
\AtBeginDocument{\ShorthandsOn}
\AtBeginEnvironment{tikzpicture}{\ShorthandsOff}

\RequirePackage{geometry}
\RequirePackage{fancyhdr}
\RequirePackage{lastpage}
\RequirePackage{multicol}
\RequirePackage{array}
\RequirePackage{enumitem}
\RequirePackage{xcolor}
\RequirePackage{needspace}

\RequirePackage{amsmath,amssymb}
\iflycee@cm\else\RequirePackage{fourier}\fi
\DeclareMathAlphabet{\mathcal}{OMS}{cmsy}{m}{n}
\renewcommand{\ttdefault}{lmtt}
\begingroup\catcode`\_=\active \catcode`\^=\active
\global\let_\sb \global\let^\sp
\endgroup
\AtBeginDocument{\catcode`\_=12 \mathcode`\_="8000
  \catcode`\^=12 \mathcode`\^="8000 }
\newlength{\retraitcalculs}
\setlength{\retraitcalculs}{2em}

\RequirePackage{tikz}
\usetikzlibrary{arrows.meta,calc,babel,shapes.misc}
\RequirePackage[most]{tcolorbox}
\tcbuselibrary{skins,breakable}

\definecolor{coulDef}{HTML}{1F4E79}
\definecolor{coulProp}{HTML}{7030A0}
\definecolor{coulRem}{HTML}{C55A11}
\definecolor{coulEx}{HTML}{2E7D32}
\definecolor{coulExo}{HTML}{B03A2E}
\definecolor{coulPart}{HTML}{1F4E79}
\definecolor{coulMeth}{HTML}{1B6E4F}
\definecolor{coulCourbe}{HTML}{0B5ED7}
\definecolor{coulCourbeB}{HTML}{C00000}
\definecolor{coulCourbeC}{HTML}{2E7D32}
\definecolor{coulGrille}{HTML}{8C8C8C}
\definecolor{coulRep}{HTML}{1B6E4F}
\definecolor{coulBar}{HTML}{2E7D32}

\newcommand{\R}{\mathbb{R}}
\newcommand{\Rs}{\mathbb{R}^{*}}
\newcommand{\N}{\mathbb{N}}
\newcommand{\Ns}{\mathbb{N}^{*}}
\newcommand{\Z}{\mathbb{Z}}
\newcommand{\Q}{\mathbb{Q}}
\renewcommand{\le}{\leqslant}
\renewcommand{\ge}{\geqslant}
\renewcommand{\approx}{\simeq}
\newcommand{\vect}[1]{\overrightarrow{#1}}
\newcommand{\norme}[1]{\left\lVert #1 \right\rVert}
\newcommand{\intervalle}[2]{\left[#1\,;#2\right]}
\RequirePackage{cancel}
\renewcommand{\CancelColor}{\color{coulExo}}
\newcommand{\fois}[1]{\times{\color{coulExo}#1}}
\newcommand{\barre}[1]{\cancel{#1}}
\newlength{\lycee@eqh}\newlength{\lycee@eqd}
\newsavebox{\lycee@eqa}\newsavebox{\lycee@eqb}\newsavebox{\lycee@eqc}
\newcommand{\lycee@eqcol}[1]{\hbox{$\renewcommand{\arraystretch}{1.35}%
  \begin{array}[t]{@{}>{\displaystyle}r@{}>{\displaystyle{}}l@{}}#1\end{array}$}}
\newcommand{\lycee@eqmesure}[1]{%
  \ifdim\ht#1>\lycee@eqh\setlength{\lycee@eqh}{\ht#1}\fi
  \ifdim\dp#1>\lycee@eqd\setlength{\lycee@eqd}{\dp#1}\fi}
\newcommand{\lycee@eqtrait}{\hspace{1.2em}ou\hspace{1.2em}}
\newcommand{\lycee@eqiff}{$\iff$\hspace{0.8em}}
\newcommand{\eqproduit}[2]{\par\smallskip\noindent C'est une équation produit :
  $A\times B=0 \iff A=0$ ou $B=0$.\par\nopagebreak\smallskip
  \begingroup\setlength{\lycee@eqh}{0pt}\setlength{\lycee@eqd}{0pt}%
  \setbox\lycee@eqa=\lycee@eqcol{#1}\setbox\lycee@eqb=\lycee@eqcol{#2}%
  \lycee@eqmesure\lycee@eqa\lycee@eqmesure\lycee@eqb%
  \noindent\hspace*{\retraitcalculs}\lycee@eqiff\usebox\lycee@eqa\lycee@eqtrait\usebox\lycee@eqb\endgroup\par\smallskip}
\newcommand{\eqproduittrois}[3]{\par\smallskip\noindent C'est une équation produit :
  $A\times B\times C=0 \iff A=0$ ou $B=0$ ou $C=0$.\par\nopagebreak\smallskip
  \begingroup\setlength{\lycee@eqh}{0pt}\setlength{\lycee@eqd}{0pt}%
  \setbox\lycee@eqa=\lycee@eqcol{#1}\setbox\lycee@eqb=\lycee@eqcol{#2}\setbox\lycee@eqc=\lycee@eqcol{#3}%
  \lycee@eqmesure\lycee@eqa\lycee@eqmesure\lycee@eqb\lycee@eqmesure\lycee@eqc%
  \noindent\hspace*{\retraitcalculs}\lycee@eqiff\usebox\lycee@eqa\lycee@eqtrait\usebox\lycee@eqb\lycee@eqtrait\usebox\lycee@eqc\endgroup
  \par\smallskip}

\tikzset{grille/.style={coulGrille,line width=0.4pt}}
\newcommand{\quadrillage}[4]{\draw[grille,step=1] (#1,#2) grid (#3,#4);}
\tikzset{
  croix/.style={cross out,draw,line width=0.6pt,inner sep=0pt,minimum size=4pt},
  croixdroite/.style={croix,rotate=45,line width=0.8pt},
}
\newcommand{\pt}[3]{\path (#1) node[croix]{} node[#2,font=\small,inner sep=2.5pt]{$#3$};}

\newcommand{\lycee@repere}[4]{%
  \draw[grille,step=1] (#1,#2) grid (#3,#4);
  \draw[-{Stealth[length=2mm]},thick] (#1,0)--({#3+0.4},0) node[below left=-1pt]{$x$};
  \draw[-{Stealth[length=2mm]},thick] (0,#2)--(0,{#4+0.4}) node[below left=-1pt]{$y$};
  \node[below left,font=\scriptsize,inner sep=1.5pt] at (0,0) {$0$};}
\newcommand{\lycee@gradx}[1]{\foreach \k/\l in {#1}{%
  \draw (\k,0.07)--(\k,-0.07) node[below,font=\scriptsize,inner sep=1.5pt]{$\l$};}}
\newcommand{\lycee@grady}[1]{\foreach \k/\l in {#1}{%
  \draw (0.07,\k)--(-0.07,\k) node[left,font=\scriptsize,inner sep=1.5pt]{$\l$};}}
\newcommand{\lycee@intff}[2]{\left[#1\,;#2\right]}
\newcommand{\lycee@intfo}[2]{\left[#1\,;#2\right[}
\newcommand{\lycee@intof}[2]{\left]#1\,;#2\right]}
\newcommand{\lycee@intoo}[2]{\left]#1\,;#2\right[}
\newcommand{\lycee@ens}[1]{\left\{#1\right\}}
\AtBeginDocument{%
  \providecommand{\repere}{\lycee@repere}%
  \providecommand{\gradx}{\lycee@gradx}%
  \providecommand{\grady}{\lycee@grady}%
  \providecommand{\Cf}{\mathcal{C}\sb{f}}%
  \providecommand{\Cg}{\mathcal{C}\sb{g}}%
  \providecommand{\intff}{\lycee@intff}%
  \providecommand{\intfo}{\lycee@intfo}%
  \providecommand{\intof}{\lycee@intof}%
  \providecommand{\intoo}{\lycee@intoo}%
  \providecommand{\ens}{\lycee@ens}}

\newlist{questions}{enumerate}{3}
\setlist[questions,1]{label=\arabic*\textdegree),leftmargin=*,itemsep=2pt,topsep=3pt}
\setlist[questions,2]{label=\alph*),leftmargin=*,itemsep=1pt,topsep=2pt}
\setlist[questions,3]{label=\roman*),leftmargin=*,itemsep=1pt,topsep=1pt}
\setlist[itemize]{label=\textbullet,leftmargin=*,itemsep=2pt,topsep=3pt}

\newcommand{\besoinplace}[1]{\par\penalty\@M\begingroup
  \dimen@=#1\relax\dimen@ii=\pagegoal\advance\dimen@ii-\pagetotal
  \ifdim\dimen@>\dimen@ii\ifdim\dimen@ii>\z@\vfil\fi\break\fi\endgroup}

\newcommand{\lycee@pied}{\thepage/\pageref{LastPage}}
\newcommand{\entete}[2]{%
  \pagestyle{fancy}\fancyhf{}%
  \fancyhead[L]{\footnotesize\color{coulPart}\textbf{#1}}%
  \fancyhead[R]{\footnotesize\color{coulPart}\textbf{#2}}%
  \fancyfoot[C]{\footnotesize\lycee@pied}%
  \renewcommand{\headrulewidth}{0.4pt}%
  \renewcommand{\headrule}{\hbox to\headwidth{%
    \color{coulPart!40}\leaders\hrule height \headrulewidth\hfill}}}

\RequirePackage{ccicons}
\newcommand{\lycee@licence}{{\scriptsize\color{gray}%
  \copyright\ Luc Destombe --- \ccbyncsa\ CC BY-NC-SA 4.0}}
\newif\iflycee@licence \lycee@licencetrue
\newcommand{\sanslicence}{\lycee@licencefalse}
\AtBeginDocument{\iflycee@licence\AddToHookNext{shipout/foreground}{%
  \put(0,0){\rlap{\kern\dimexpr1in+\hoffset+\oddsidemargin+\textwidth\relax
    \llap{\raisebox{-\dimexpr1in+\voffset+\topmargin+\headheight+\headsep
      +\textheight+\footskip\relax}{\lycee@licence}}}}}\fi}

\newcommand{\lycee@site}{faire-des-maths.fr}
\newcommand{\lycee@version}{{\scriptsize\color{gray}\lycee@site
  \ifdefined\lyceeversion\ --- Version~\lyceeversion\fi}}
\AtBeginDocument{\iflycee@licence\AddToHookNext{shipout/foreground}{%
  \put(0,0){\rlap{\kern\dimexpr1in+\hoffset+\oddsidemargin\relax
    \raisebox{-\dimexpr1in+\voffset+\topmargin+\headheight+\headsep
      +\textheight+\footskip\relax}{\lycee@version}}}}\fi}

\newcommand{\formulesserrees}{%
  \setlength{\abovedisplayskip}{3pt}\setlength{\belowdisplayskip}{3pt}%
  \setlength{\abovedisplayshortskip}{2pt}\setlength{\belowdisplayshortskip}{3pt}}

\setlength{\parindent}{0pt}

\RequirePackage[implicit=false,hidelinks,pdfencoding=auto,psdextra,bookmarksopen,
  bookmarksopenlevel=1,pdfcreator={},pdfproducer={}]{hyperref}
\RequirePackage{bookmark}
\hypersetup{pdfauthor={Luc Destombe},pdfkeywords={CC BY-NC-SA 4.0}}
\newcounter{lycee@signet}
\newcommand{\signet}[3][]{\stepcounter{lycee@signet}%
  \edef\lycee@dest{\if\relax\detokenize{#1}\relax
    signet.\arabic{lycee@signet}\else#1\fi}%
  \hypertarget{\lycee@dest}{}\bookmark[level=#2,dest=\lycee@dest]{#3}}
\pdfstringdefDisableCommands{%
  \def\R{R}\def\N{N}\def\Z{Z}\def\Q{Q}\def\vect#1{#1}\def\Cf{Cf}\def\Cg{Cg}\def\fois#1{×#1}%
  \def\barre#1{#1}\def\intervalle#1#2{[#1;#2]}\def\intff#1#2{[#1;#2]}%
  \def\textsuperscript#1{#1}\def\dfrac#1#2{#1/#2}\def\frac#1#2{#1/#2}%
  \def\vec#1{#1⃗}}%
\newcommand{\res}[1]{#1}
\geometry{top=2cm,bottom=1cm,left=1cm,right=1cm,headheight=14pt,footskip=0.6cm}

\RequirePackage{pgfplots}
\pgfplotsset{compat=1.18}
\RequirePackage{tkz-tab}

\newcounter{partiecours}
\renewcommand{\thepartiecours}{\Roman{partiecours}}
\newcommand{\partiecours}[1]{%
  \refstepcounter{partiecours}%
  \Needspace*{8\baselineskip}%
  \bigskip
  \noindent\begin{tcolorbox}[enhanced,colback=coulPart,colframe=coulPart,
      boxrule=0pt,arc=2pt,left=8pt,right=8pt,top=4pt,bottom=4pt,
      before skip=6pt,after skip=10pt]
    \color{white}\bfseries\large \thepartiecours{} --- #1%
    \signet[partie-\arabic{partiecours}]{1}{\thepartiecours{} --- #1}
  \end{tcolorbox}%
  \addcontentsline{toc}{section}{\thepartiecours{} --- #1}%
}

\newtcolorbox{boitecours}[3][]{%
  enhanced,breakable,
  colback=white,colframe=#2,coltitle=white,
  fonttitle=\bfseries,
  attach boxed title to top left={xshift=8pt,yshift=-\tcboxedtitleheight/2},
  boxed title style={colback=#2,arc=2pt,boxrule=0pt},
  boxrule=0.9pt,arc=2pt,
  left=8pt,right=8pt,top=8pt,bottom=6pt,
  title={#3},#1}

\newcounter{methode}
\newenvironment{definitionbox}[1][Définition]%
  {\begin{boitecours}[phantom={\signet{2}{#1}}]{coulDef}{#1}}{\end{boitecours}}
\newenvironment{proprietebox}[1][Propriété]%
  {\begin{boitecours}[phantom={\signet{2}{#1}}]{coulProp}{#1}}{\end{boitecours}}
\newenvironment{methodebox}[1][Méthode]{\stepcounter{methode}%
  \begin{boitecours}[phantom={\signet[methode-\arabic{methode}]{2}{#1}}]{coulMeth}{#1}}%
  {\end{boitecours}}

\newcommand{\etape}[1]{\par\smallskip\textbf{\color{coulMeth}#1}\par\nobreak\smallskip}
\newcommand{\enonce}[1]{\emph{#1}\par\smallskip}
\newcommand{\reponse}[1]{\par\smallskip
  {\footnotesize\itshape\color{black!55}Réponses : #1}}
\newcommand{\demo}[1]{\textbf{\color{coulMeth}Démonstration.} #1}
\newcommand{\afaire}[1]{%
  \par\smallskip
  \noindent{\small\color{coulExo}$\blacktriangleright$\ \textbf{À faire :}
    fiche d'exercices du chapitre, #1.}\par\smallskip}

\newenvironment{calculs}{\setlength{\jot}{5pt}\@fleqntrue\@mathmargin\retraitcalculs
  \setlength{\abovedisplayskip}{4pt}\setlength{\belowdisplayskip}{4pt}%
  \setlength{\abovedisplayshortskip}{4pt}\setlength{\belowdisplayshortskip}{4pt}%
  \csname alignat*\endcsname{2}}%
  {\csname endalignat*\endcsname}
\newcommand{\rem}[1]{\text{\small\itshape\color{black!60}#1}}
\newcommand{\explique}[2]{\par\smallskip\noindent
  \begin{minipage}[t]{0.57\linewidth}\raggedright #1\end{minipage}\hfill
  \begin{minipage}[t]{0.40\linewidth}\raggedright\leavevmode
    {\small\itshape\color{black!60}#2}\end{minipage}\par}
\newcommand{\cas}[1]{\par\smallskip\noindent\textbf{\color{coulExo}#1}\vadjust{\nobreak}}
\newcommand{\calc}[1]{\par\smallskip\textbf{\color{coulExo}#1}\par\nobreak}

\newtcolorbox{remarquebox}[1][Remarque]{%
  enhanced,breakable,colback=white,colframe=coulRem,
  boxrule=0pt,leftrule=2.5pt,arc=0pt,outer arc=0pt,
  left=8pt,right=6pt,top=5pt,bottom=5pt,
  before upper={\textbf{\color{coulRem}#1}\par\smallskip}}

\newtcolorbox{exemplebox}[1][Exemple]{%
  enhanced,breakable,colback=white,colframe=coulEx,
  boxrule=0pt,leftrule=2.5pt,arc=0pt,outer arc=0pt,
  left=8pt,right=6pt,top=5pt,bottom=5pt,
  before upper={\textbf{\color{coulEx}#1}\par\smallskip}}

\pgfplotsset{
  repere/.style={
    axis lines=middle,
    axis line style={-{Stealth[length=2.4mm]},thick},
    every axis x label/.style={at={(ticklabel* cs:1.0)},anchor=west},
    every axis y label/.style={at={(ticklabel* cs:1.0)},anchor=south},
    label style={font=\footnotesize},
    tick label style={font=\scriptsize},
    grid=both,
    major grid style={grille},
    minor tick num=0,
    tick align=inside,
    clip mode=individual,
  },
}
\tikzset{
  courbe/.style={coulCourbe,line width=1.1pt,smooth},
  courbeB/.style={coulCourbeB,line width=1.1pt,smooth},
  courbeC/.style={coulCourbeC,line width=1.1pt,smooth},
}

\newlist{etapes}{enumerate}{1}
\setlist[etapes]{label=\textbf{\arabic*.},leftmargin=*,itemsep=1pt,topsep=2pt}

\newlist{expressions}{enumerate}{1}
\setlist[expressions]{label={\Alph*\ $=$},leftmargin=*,itemsep=3pt,topsep=3pt}

\newcounter{exo}
\renewcommand{\theexo}{\ifnum\value{exo}<10 0\fi\arabic{exo}}
\newtcolorbox{exercice}[1][]{%
  enhanced,breakable,
  colback=white,colframe=coulExo,
  boxrule=0.9pt,arc=2pt,
  left=8pt,right=8pt,top=8pt,bottom=6pt,
  coltitle=white,fonttitle=\bfseries,
  attach boxed title to top left={xshift=8pt,yshift=-\tcboxedtitleheight/2},
  boxed title style={colback=coulExo,arc=2pt,boxrule=0pt},
  phantom={\signet{2}{Exercice \arabic{exo}}},
  title={Exercice \theexo\ifx\relax#1\relax\else{} --- #1\fi}}
\newenvironment{exo}[1][\relax]{\refstepcounter{exo}\begin{exercice}[#1]}{\end{exercice}}

  \newenvironment{correction}{\par}{\par}
  \newcommand{\autableau}{}
  \newcommand{\etapeexemples}{\etape{Exemples corrigés}}
  \newcommand{\etapetableau}[1]{\etape{#1}}
  \newenvironment{exempletableau}[1][Exemple]%
    {\begin{exemplebox}[#1]}{\end{exemplebox}}

\RequirePackage{comment}
\excludecomment{correctionavj}

\newcommand{\coursEntete}{}
\newcommand{\coursNiveau}{}
\newcommand{\coursTitre}{}
\newcommand{\coursSurTitre}{}
\newcommand{\coursSousTitre}{}

\newcommand{\enteteCours}[1]{\renewcommand{\coursEntete}{#1}}
\newcommand{\chapitrecours}[2]{\enteteCours{Chapitre #1 --- #2}}
\newcommand{\niveaucours}[1]{\renewcommand{\coursNiveau}{#1}}
\newcommand{\titrecours}[3][]{%
  \renewcommand{\coursTitre}{#2}%
  \renewcommand{\coursSurTitre}{#1}%
  \renewcommand{\coursSousTitre}{#3}}

\pagestyle{fancy}
\fancyhf{}
\fancyhead[L]{\footnotesize\color{coulPart}\textbf{\coursEntete}}
\fancyhead[R]{\footnotesize\color{coulPart}\textbf{\coursNiveau}}
\fancyfoot[C]{\footnotesize\thepage}
\renewcommand{\headrulewidth}{0.4pt}
\renewcommand{\headrule}{\hbox to\headwidth{\color{coulPart!40}\leaders\hrule height \headrulewidth\hfill}}

\setlength{\parskip}{2pt}

\newcommand{\lycee@mentioneleve}{\\[2pt]{\small\itshape Version élève}}
\AtBeginDocument{%
  \begin{center}
    \begin{tcolorbox}[enhanced,width=0.72\linewidth,colback=white,colframe=coulPart,
        boxrule=1.6pt,arc=3pt,halign=center,top=10pt,bottom=10pt]
      \ifx\coursSurTitre\empty
        {\Huge\bfseries\color{coulPart} \coursTitre}\\[4pt]
      \else
        {\Huge\bfseries\color{coulPart} \coursTitre}\\[3pt]
        {\Large\bfseries\color{coulPart} \coursSurTitre}\\[5pt]
      \fi
      {\normalsize \coursSousTitre}%
      
    \end{tcolorbox}
  \end{center}%
}

\RequirePackage{tikz}
\RequirePackage{xcolor}

\usetikzlibrary{decorations.text}
\definecolor{lycee@metal}{HTML}{1F4E79}
\definecolor{lycee@manche}{HTML}{8A5A2B}
\definecolor{lycee@etincelle}{HTML}{F28C28}
\newcommand{\lycee@pluriel}[2]{#1 #2\ifnum#1>1 s\fi}
\newcommand{\lycee@sceau}{%
  \edef\lycee@bas{\lycee@pluriel{\lyceerelectures}{relecture} et
    \lycee@pluriel{\lyceerodages}{rodage}}%
  \begin{tikzpicture}[line join=round]
    \fill[coulPart!10,draw=coulPart,line width=.8pt] (0:2.2)
      \foreach \i in {1,...,71}{ -- (\i*5:{isodd(\i)?2.1:2.2}) } -- cycle;
    \draw[coulPart,line width=1.2pt] (0,0) circle (1.93);
    \draw[coulPart,line width=.4pt] (0,0) circle (1.85);
    \draw[coulPart,line width=.4pt] (0,0) circle (1.2);
    \path[decorate,decoration={text along path,text align=center,
        text={|\fontsize{9.5}{11}\selectfont\bfseries\color{coulPart}|TEST{\'E} AVEC MES {\'E}L{\`E}VES}}]
      (180:1.45) arc (180:0:1.45);
    \path[decorate,decoration={text along path,text align=center,
        text/.expanded={|\noexpand\fontsize{8.5}{10}\noexpand\selectfont\noexpand\color{coulPart}|\lycee@bas}}]
      (195:1.68) arc (195:345:1.68);
    \begin{scope}[x=.25mm,y=-.25mm]
      \begin{scope}[shift={(-32,-40)}]%
        \fill[lycee@metal] (23,31) rectangle (41,49);
        \draw[lycee@etincelle,line width=.5pt,line cap=round]
          (20,30)--(15,27) (19,34)--(14,35) (44,30)--(49,27) (45,34)--(50,35);
        \fill[lycee@manche,rounded corners=.3mm] (38,24.6) rectangle (60,27.6);
        \fill[lycee@metal,rounded corners=.4mm] (25,20) rectangle (39,31);
        \fill[lycee@etincelle,rounded corners=.2mm] (25,20) rectangle (28,31);
      \end{scope}
    \end{scope}
    \ifdefined\lyceesceauplusieurs\def\lycee@annee{Ann{\'e}es~:}\else
      \def\lycee@annee{Ann{\'e}e}\fi
    \path[decorate,decoration={text along path,text align=center,
        text/.expanded={|\noexpand\fontsize{8}{10}\noexpand\selectfont\noexpand\color{coulPart}|%
          \unexpanded\expandafter{\lycee@annee} \unexpanded\expandafter{\lyceesceauannees}}}]
      (95:1.08) arc (95:445:1.08);
  \end{tikzpicture}}
\AtEndDocument{\ifdefined\lyceesceauannees
  \par\besoinplace{5cm}\vspace{1.5em}%
  \begin{center}\lycee@sceau\end{center}\fi}
\makeatother

\chapitrecours{8}{Orthogonalité et distances}
\niveaucours{T\textsuperscript{le} spécialité}
\titrecours{ORTHOGONALITÉ ET DISTANCES}{Cours --- Terminale spécialité}

\begin{document}

\newcommand{\cubeespace}[1]{%
  \coordinate (A) at (0,0); \coordinate (B) at (#1,0);
  \coordinate (D) at ({0.45*#1},{0.35*#1}); \coordinate (C) at ({#1+0.45*#1},{0.35*#1});
  \coordinate (E) at (0,#1); \coordinate (F) at (#1,#1);
  \coordinate (H) at ({0.45*#1},{#1+0.35*#1}); \coordinate (G) at ({#1+0.45*#1},{#1+0.35*#1});
  \draw (A)--(B)--(F)--(E)--cycle (B)--(C)--(G)--(F) (E)--(H)--(G);
  \draw[dashed] (A)--(D)--(C) (D)--(H);
  \path[font=\footnotesize] (A) node[below left]{$A$} (B) node[below right]{$B$}
    (C) node[right]{$C$} (D) node[above left]{$D$} (E) node[left]{$E$}
    (F) node[below right]{$F$} (G) node[above right]{$G$} (H) node[above left]{$H$};}
\newcommand{\accol}[1]{\left\{\renewcommand{\arraystretch}{1.3}\begin{array}{rcl}#1\end{array}\right.}

\begin{remarquebox}[Durée du chapitre]
À titre indicatif, ce chapitre est prévu pour un volume de $12$~h d'enseignement
($6$ séances de $2$~h). Une partie importante des exercices est à faire en autonomie.
\end{remarquebox}

\partiecours{Rappels de 1re : le produit scalaire dans le plan}

\begin{minipage}[c]{0.62\linewidth}\raggedright
\begin{definitionbox}[Définition 1 --- Produit scalaire dans le plan]
Soit $\vect{u}=\vect{AB}$ et $\vect{v}=\vect{AC}$ deux vecteurs non nuls du plan. Le
\textbf{produit scalaire} de $\vect{u}$ et $\vect{v}$ est le nombre réel
\[
\vect{u}\cdot\vect{v}=AB\times AC\times\cos\widehat{BAC}.
\]
Si l'un des deux vecteurs est nul, $\vect{u}\cdot\vect{v}=0$.
\end{definitionbox}
\end{minipage}\hfill
\begin{minipage}[c]{0.34\linewidth}\centering
\begin{tikzpicture}[font=\footnotesize]
  \coordinate (A) at (0,0); \coordinate (B) at (3,0); \coordinate (C) at (2,1.6);
  \coordinate (K) at (2,0);
  \draw[-{Stealth},coulExo,thick] (A)--(B) node[midway,below]{$\vect{u}$};
  \draw[-{Stealth},coulCourbe,thick] (A)--(C) node[midway,above left]{$\vect{v}$};
  \draw[dashed] (C)--(K);
  \draw (K)++(-0.2,0)--++(0,0.2)--++(0.2,0);
  \path (A) node[below left]{$A$} (B) node[below right]{$B$} (C) node[above]{$C$}
        (K) node[below]{$K$};
\end{tikzpicture}
\end{minipage}

\begin{proprietebox}[Propriété 1 --- Autres expressions, propriétés]
\begin{itemize}
  \item \textbf{Projeté orthogonal.} Si $K$ est le projeté orthogonal de $C$ sur la
        droite $(AB)$ : $\vect{AB}\cdot\vect{AC}=\vect{AB}\cdot\vect{AK}$, qui vaut
        $AB\times AK$ si $\vect{AB}$ et $\vect{AK}$ sont de même sens, $-AB\times AK$
        sinon.
  \item \textbf{Normes.} $\vect{u}\cdot\vect{u}=\|\vect{u}\|^{2}$ (carré scalaire) et,
        en développant :
        \begin{center}
        $\|\vect{u}+\vect{v}\|^{2}=\|\vect{u}\|^{2}+2\,\vect{u}\cdot\vect{v}+\|\vect{v}\|^{2}$ \qquad
        $\|\vect{u}-\vect{v}\|^{2}=\|\vect{u}\|^{2}-2\,\vect{u}\cdot\vect{v}+\|\vect{v}\|^{2}$
        \end{center}
  \item \textbf{Formules de polarisation} (on isole $\vect{u}\cdot\vect{v}$ dans les
        égalités précédentes) :
        \begin{center}
        $\vect{u}\cdot\vect{v}=\frac{1}{2}\left(\|\vect{u}+\vect{v}\|^{2}-\|\vect{u}\|^{2}-\|\vect{v}\|^{2}\right)$ \qquad
        $\vect{u}\cdot\vect{v}=\frac{1}{2}\left(\|\vect{u}\|^{2}+\|\vect{v}\|^{2}-\|\vect{u}-\vect{v}\|^{2}\right)$
        \end{center}
  \item \textbf{Coordonnées.} Dans un repère orthonormé du plan, si
        $\vect{u}\begin{pmatrix}x\\y\end{pmatrix}$ et $\vect{v}\begin{pmatrix}x'\\y'\end{pmatrix}$, alors
        $\vect{u}\cdot\vect{v}=xx'+yy'$.
  \item \textbf{Calcul.} Pour tous vecteurs $\vect{u}$, $\vect{v}$, $\vect{w}$ et tout
        nombre réel $k$ : $\vect{u}\cdot\vect{v}=\vect{v}\cdot\vect{u}$ (symétrie),
        $\vect{u}\cdot(\vect{v}+\vect{w})=\vect{u}\cdot\vect{v}+\vect{u}\cdot\vect{w}$ et
        $\vect{u}\cdot(k\vect{v})=k\,(\vect{u}\cdot\vect{v})$ (bilinéarité).
  \item \textbf{Orthogonalité.} $\vect{u}$ et $\vect{v}$ sont orthogonaux si et
        seulement si $\vect{u}\cdot\vect{v}=0$.
\end{itemize}
\end{proprietebox}

\begin{methodebox}[Méthode 1 --- Calculer un produit scalaire en choisissant la bonne expression]
\etapeexemples
Calculer chaque produit scalaire.
\begin{questions}[label=\alph*)]
  \item $\vect{AB}\cdot\vect{AC}$, dans un triangle $ABC$ où $AB=4$, $AC=5$ et
        $\widehat{BAC}=\frac{\pi}{3}$.
  \item $\vect{AB}\cdot\vect{AI}$, où $ABCD$ est un carré de côté $3$ et $I$ le milieu
        de $[CD]$.
  \item $\vect{u}\cdot\vect{v}$, avec $\vect{u}\begin{pmatrix}3\\-1\end{pmatrix}$ et
        $\vect{v}\begin{pmatrix}2\\4\end{pmatrix}$ dans un repère orthonormé.
  \item $\vect{u}\cdot\vect{v}$, sachant que $\|\vect{u}\|=3$, $\|\vect{v}\|=4$ et
        $\|\vect{u}+\vect{v}\|=6$.
\end{questions}
\begin{correction}
\cas{a)}
\begin{calculs}
\vect{AB}\cdot\vect{AC} &= AB\times AC\times\cos\widehat{BAC} &\qquad& \rem{on connaît l'angle}\\
\vect{AB}\cdot\vect{AC} &= 4\times 5\times\frac{1}{2} && \rem{$\cos\left(\frac{\pi}{3}\right)=\frac{1}{2}$}\\
\vect{AB}\cdot\vect{AC} &= \res{10} &&
\end{calculs}
\cas{b)}
\explique{Le projeté orthogonal de $I$ sur $(AB)$ est le milieu de $[AB]$ : il est à
la distance $\frac{3}{2}$ de $A$, du côté de $B$.}{$(AB)$ et $(CD)$ parallèles}
\begin{calculs}
\vect{AB}\cdot\vect{AI} &= AB\times\frac{AB}{2} &\qquad& \rem{projeté, même sens}\\
\vect{AB}\cdot\vect{AI} &= 3\times\frac{3}{2} &&\\
\vect{AB}\cdot\vect{AI} &= \res{\frac{9}{2}} &&
\end{calculs}
\cas{c)}
\begin{calculs}
\vect{u}\cdot\vect{v} &= 3\times 2+(-1)\times 4 &\qquad& \rem{$xx'+yy'$, repère orthonormé}\\
\vect{u}\cdot\vect{v} &= \res{2} &&
\end{calculs}
\cas{d)}
\begin{calculs}
\vect{u}\cdot\vect{v} &= \frac{1}{2}\left(\|\vect{u}+\vect{v}\|^{2}-\|\vect{u}\|^{2}-\|\vect{v}\|^{2}\right) &\qquad& \rem{polarisation}\\
\vect{u}\cdot\vect{v} &= \frac{1}{2}\left(36-9-16\right) &&\\
\vect{u}\cdot\vect{v} &= \res{\frac{11}{2}} &&
\end{calculs}
\end{correction}

\etape{À vous de jouer}
Calculer chaque produit scalaire.
\begin{questions}[label=\alph*)]
  \item $\vect{AB}\cdot\vect{AC}$, dans un triangle $ABC$ où $AB=6$, $AC=2$ et
        $\widehat{BAC}=\frac{\pi}{4}$.
  \item $\vect{AB}\cdot\vect{AJ}$, où $ABCD$ est un carré de côté $4$ et $J$ le milieu
        de $[BC]$.
  \item $\vect{u}\cdot\vect{v}$, avec $\vect{u}\begin{pmatrix}-2\\5\end{pmatrix}$ et
        $\vect{v}\begin{pmatrix}4\\1\end{pmatrix}$ dans un repère orthonormé.
  \item $\vect{u}\cdot\vect{v}$, sachant que $\|\vect{u}\|=2$, $\|\vect{v}\|=5$ et
        $\|\vect{u}+\vect{v}\|=4$.
\end{questions}
\reponse{a)~$6\sqrt{2}$ ;\; b)~$16$ ;\; c)~$-3$ ;\; d)~$-\frac{13}{2}$.}
\end{methodebox}

\afaire{exercices 1 à 4}

\partiecours{Produit scalaire dans l'espace}

\begin{minipage}[c]{0.58\linewidth}\raggedright
Deux vecteurs de l'espace, placés à partir d'un même point, $\vect{u}=\vect{AB}$ et
$\vect{v}=\vect{AC}$, sont toujours dans un même plan : un plan qui contient $A$, $B$
et $C$. On y calcule leur produit scalaire comme en 1re.
\end{minipage}\hfill
\begin{minipage}[c]{0.38\linewidth}\centering
\begin{tikzpicture}[font=\footnotesize]
  \cubeespace{2.2}
  \fill[coulExo,opacity=0.12] (A)--(F)--(H)--cycle;
  \draw[-{Stealth},coulExo,thick] (A)--(F) node[midway,below right]{$\vect{u}$};
  \draw[-{Stealth},coulCourbe,thick] (A)--(H) node[midway,above left]{$\vect{v}$};
\end{tikzpicture}

{\footnotesize Exemple dans un cube : $\vect{u}=\vect{AF}$ et $\vect{v}=\vect{AH}$ sont dans le plan $(AFH)$.}
\end{minipage}

\begin{definitionbox}[Définition 2 --- Produit scalaire dans l'espace]
Soit $\vect{u}$ et $\vect{v}$ deux vecteurs de l'espace, et $A$, $B$, $C$ trois points
tels que $\vect{u}=\vect{AB}$ et $\vect{v}=\vect{AC}$. Le \textbf{produit scalaire}
$\vect{u}\cdot\vect{v}$ est le produit scalaire $\vect{AB}\cdot\vect{AC}$ calculé dans
un plan contenant $A$, $B$ et $C$ :
\[
\vect{u}\cdot\vect{v}=AB\times AC\times\cos\widehat{BAC}\qquad(\text{et }0\text{ si l'un des vecteurs est nul}).
\]
$\vect{u}$ et $\vect{v}$ sont \textbf{orthogonaux} si l'un d'eux est nul, ou si les
droites $(AB)$ et $(AC)$ sont perpendiculaires.
\end{definitionbox}

\begin{proprietebox}[Propriété 2 --- Les règles du plan restent vraies]
Toutes les expressions et règles de la propriété 1 (sauf celle des coordonnées, qui
vient en partie III) restent vraies dans l'espace : projeté orthogonal, normes et
formules de polarisation, symétrie et bilinéarité. En particulier, $\vect{u}$ et
$\vect{v}$ sont orthogonaux si et seulement si $\vect{u}\cdot\vect{v}=0$.
\end{proprietebox}

\begin{methodebox}[Méthode 2 --- Calculer un produit scalaire dans un solide]
\etapeexemples
\begin{minipage}[t]{0.6\linewidth}\raggedright
$ABCDEFGH$ est un cube d'arête $3$. Calculer :
\begin{questions}[label=\alph*)]
  \item $\vect{AB}\cdot\vect{AF}$ ;
  \item $\vect{AB}\cdot\vect{GH}$ ;
  \item $\vect{AF}\cdot\vect{AH}$ ;
  \item $\vect{AE}\cdot\vect{BH}$.
\end{questions}
\end{minipage}\hfill
\begin{minipage}[t]{0.36\linewidth}\centering\vspace{0pt}
\begin{tikzpicture}
  \cubeespace{2}
\end{tikzpicture}
\end{minipage}
\begin{correction}
\cas{a)}\par\nopagebreak
\explique{Dans le carré $ABFE$, le projeté orthogonal de $F$ sur $(AB)$ est $B$.}{on
travaille dans une face}
\begin{minipage}[c]{0.62\linewidth}\raggedright
\begin{calculs}
\vect{AB}\cdot\vect{AF} &= \vect{AB}\cdot\vect{AB} &\qquad& \rem{projeté orthogonal}\\
\vect{AB}\cdot\vect{AF} &= AB^{2} &&\\
\vect{AB}\cdot\vect{AF} &= \res{9} &&
\end{calculs}
\end{minipage}\hfill
\begin{minipage}[c]{0.34\linewidth}\centering
\begin{tikzpicture}[font=\footnotesize]
  \cubeespace{2}
  \fill[coulExo,opacity=0.12] (A)--(B)--(F)--(E)--cycle;
  \draw[-{Stealth},coulExo,thick] (A)--(B) node[midway,below]{$\vect{AB}$};
  \draw[-{Stealth},coulCourbe,thick] (A)--(F) node[pos=0.6,below right]{$\vect{AF}$};
  \draw ($(B)+(-0.2,0)$)--++(0,0.2)--++(0.2,0);
\end{tikzpicture}
\end{minipage}
\cas{b)}\par\nopagebreak
\explique{$\vect{GH}=\vect{CD}=\vect{BA}$ : on remplace par un vecteur égal, de
sens contraire à $\vect{AB}$.}{$ABCD$ et $CDHG$ sont des carrés}
\begin{minipage}[c]{0.62\linewidth}\raggedright
\begin{calculs}
\vect{AB}\cdot\vect{GH} &= \vect{AB}\cdot\vect{BA} &\qquad& \rem{$\vect{BA}=-\vect{AB}$}\\
\vect{AB}\cdot\vect{GH} &= -AB^{2} &&\\
\vect{AB}\cdot\vect{GH} &= \res{-9} &&
\end{calculs}
\end{minipage}\hfill
\begin{minipage}[c]{0.34\linewidth}\centering
\begin{tikzpicture}[font=\footnotesize]
  \cubeespace{2}
  \draw[-{Stealth},coulExo,thick] (A)--(B) node[midway,below]{$\vect{AB}$};
  \draw[-{Stealth},coulCourbe,thick] (G)--(H) node[midway,above]{$\vect{GH}$};
\end{tikzpicture}
\end{minipage}
\cas{c)}\par\nopagebreak
\explique{$AF$, $AH$ et $FH$ sont des diagonales de faces, de même longueur
$3\sqrt{2}$ : le triangle $AFH$ est équilatéral, donc
$\widehat{FAH}=\frac{\pi}{3}$.}{diagonale d'un carré de côté $3$}
\begin{minipage}[c]{0.62\linewidth}\raggedright
\begin{calculs}
\vect{AF}\cdot\vect{AH} &= AF\times AH\times\cos\widehat{FAH} &\qquad& \rem{avec l'angle}\\
\vect{AF}\cdot\vect{AH} &= 3\sqrt{2}\times 3\sqrt{2}\times\frac{1}{2} &&\\
\vect{AF}\cdot\vect{AH} &= \res{9} &&
\end{calculs}
\end{minipage}\hfill
\begin{minipage}[c]{0.34\linewidth}\centering
\begin{tikzpicture}[font=\footnotesize]
  \cubeespace{2}
  \fill[coulExo,opacity=0.12] (A)--(F)--(H)--cycle;
  \draw (F)--(H);
  \draw[-{Stealth},coulExo,thick] (A)--(F) node[pos=0.6,below right]{$\vect{AF}$};
  \draw[-{Stealth},coulCourbe,thick] (A)--(H) node[midway,above left]{$\vect{AH}$};
\end{tikzpicture}
\end{minipage}
\cas{d)}\par\nopagebreak
\explique{On décompose $\vect{BH}$ par la relation de Chasles en suivant des arêtes :
$\vect{BH}=\vect{BA}+\vect{AE}+\vect{EH}$ (en pointillés sur la figure).}{des vecteurs
simples à multiplier par $\vect{AE}$}
\begin{minipage}[c]{0.62\linewidth}\raggedright
\begin{calculs}
\vect{AE}\cdot\vect{BH} &= \vect{AE}\cdot\left(\vect{BA}+\vect{AE}+\vect{EH}\right) &\qquad& \rem{Chasles}\\
\vect{AE}\cdot\vect{BH} &= \vect{AE}\cdot\vect{BA}+\vect{AE}\cdot\vect{AE}+\vect{AE}\cdot\vect{EH} && \rem{bilinéarité}\\
\vect{AE}\cdot\vect{BH} &= 0+AE^{2}+0 && \rem{arêtes perpendiculaires}\\
\vect{AE}\cdot\vect{BH} &= \res{9} &&
\end{calculs}
\end{minipage}\hfill
\begin{minipage}[c]{0.34\linewidth}\centering
\begin{tikzpicture}[font=\footnotesize]
  \cubeespace{2}
  \draw[-{Stealth},coulCourbe,thick] (B)--(H) node[pos=0.7,right]{$\vect{BH}$};
  \draw[-{Stealth},coulExo,thick] (A)--(E) node[midway,left]{$\vect{AE}$};
  \draw[coulCourbe,thick,dash pattern=on 2pt off 1.5pt] (B)--(A) (E)--(H);
\end{tikzpicture}
\end{minipage}
\end{correction}

\etape{À vous de jouer}
$ABCDEFGH$ est un cube d'arête $2$. Calculer :
\begin{questions}[label=\alph*)]
  \item $\vect{AD}\cdot\vect{AH}$ ;
  \item $\vect{AD}\cdot\vect{GF}$ ;
  \item $\vect{DE}\cdot\vect{DG}$ ;
  \item $\vect{AB}\cdot\vect{DF}$.
\end{questions}
\reponse{a)~$4$ ;\; b)~$-4$ ;\; c)~$4$ ;\; d)~$4$.}
\end{methodebox}

\afaire{exercices 5 à 9}

\partiecours{Produit scalaire dans un repère orthonormé}

\begin{definitionbox}[Définition 3 --- Base et repère orthonormés]
Une base $(\vect{i},\vect{j},\vect{k})$ de l'espace est \textbf{orthonormée} si ses
trois vecteurs sont orthogonaux deux à deux et de norme $1$. Un repère
$(O\,;\vect{i},\vect{j},\vect{k})$ est \textbf{orthonormé} si sa base l'est.

Exemple : si $ABCDEFGH$ est un cube d'arête $1$, le repère
$(A\,;\vect{AB},\vect{AD},\vect{AE})$ est orthonormé.
\end{definitionbox}

\begin{proprietebox}[Propriété 3 --- Calculs en repère orthonormé]
Dans un repère orthonormé, soit $\vect{u}\begin{pmatrix}x\\y\\z\end{pmatrix}$ et
$\vect{v}\begin{pmatrix}x'\\y'\\z'\end{pmatrix}$, et deux points $A(x_{A}\,;y_{A}\,;z_{A})$ et $B(x_{B}\,;y_{B}\,;z_{B})$.
\[
\vect{u}\cdot\vect{v}=xx'+yy'+zz'
\qquad
\|\vect{u}\|=\sqrt{x^{2}+y^{2}+z^{2}}
\qquad
AB=\sqrt{(x_{B}-x_{A})^{2}+(y_{B}-y_{A})^{2}+(z_{B}-z_{A})^{2}}
\]
\end{proprietebox}

\textbf{D'où vient la première formule ?} On écrit $\vect{u}=x\vect{i}+y\vect{j}+z\vect{k}$
et $\vect{v}=x'\vect{i}+y'\vect{j}+z'\vect{k}$, puis on développe par bilinéarité : les
neuf produits se réduisent à trois, car $\vect{i}\cdot\vect{i}$, $\vect{j}\cdot\vect{j}$ et $\vect{k}\cdot\vect{k}$ valent $1$
(norme $1$) et $\vect{i}\cdot\vect{j}$, $\vect{i}\cdot\vect{k}$ et $\vect{j}\cdot\vect{k}$ sont nuls (orthogonaux).
La norme vient de $\|\vect{u}\|^{2}=\vect{u}\cdot\vect{u}=x^{2}+y^{2}+z^{2}$.

\begin{methodebox}[Méthode 3 --- Calculer des longueurs, un produit scalaire et un angle avec des coordonnées]
\etapeexemples
Dans un repère orthonormé de l'espace, on considère les points $A(1\,;0\,;2)$,
$B(3\,;1\,;0)$ et $C(1\,;3\,;-1)$.
\begin{questions}[label=\alph*)]
  \item Calculer les longueurs $AB$ et $AC$.
  \item Calculer le produit scalaire $\vect{AB}\cdot\vect{AC}$.
  \item En déduire une mesure de l'angle $\widehat{BAC}$.
\end{questions}
\begin{correction}
\cas{a)}
\explique{On a $\vect{AB}\begin{pmatrix}3-1\\1-0\\0-2\end{pmatrix}$, soit $\vect{AB}\begin{pmatrix}2\\1\\-2\end{pmatrix}$, et
$\vect{AC}\begin{pmatrix}0\\3\\-3\end{pmatrix}$.}{coordonnées de $B$ moins celles de $A$}
\explique{Le repère est orthonormé, donc :}{condition pour utiliser les formules}
\begin{calculs}
AB &= \sqrt{2^{2}+1^{2}+(-2)^{2}} &\qquad& \\
AB &= \sqrt{9} &&\\
AB &= \res{3} &&
\end{calculs}
\begin{calculs}
AC &= \sqrt{0^{2}+3^{2}+(-3)^{2}} &\qquad& \\
AC &= \sqrt{18} &&\\
AC &= \res{3\sqrt{2}} && \rem{$\sqrt{18}=\sqrt{9\times 2}$}
\end{calculs}
\cas{b)}
\begin{calculs}
\vect{AB}\cdot\vect{AC} &= 2\times 0+1\times 3+(-2)\times(-3) &\qquad& \rem{$xx'+yy'+zz'$}\\
\vect{AB}\cdot\vect{AC} &= \res{9} &&
\end{calculs}
\cas{c)}
\explique{On a aussi $\vect{AB}\cdot\vect{AC}=AB\times AC\times\cos\widehat{BAC}$, donc
$9=3\times 3\sqrt{2}\times\cos\widehat{BAC}$ :}{les deux expressions du produit scalaire}
\begin{calculs}
\cos\widehat{BAC} &= \frac{9}{3\times 3\sqrt{2}} &\qquad& \\
\cos\widehat{BAC} &= \frac{1}{\sqrt{2}} &&\\
\cos\widehat{BAC} &= \frac{\sqrt{2}}{2} &&
\end{calculs}
\explique{Une mesure de l'angle $\widehat{BAC}$ est donc \res{$\frac{\pi}{4}$}
radian, soit $45^{\circ}$.}{un angle géométrique est entre $0$ et $\pi$}
\end{correction}

\etape{À vous de jouer}
Dans un repère orthonormé de l'espace, on considère les points $A(2\,;-1\,;0)$,
$B(3\,;0\,;0)$ et $C(2\,;0\,;1)$.
\begin{questions}[label=\alph*)]
  \item Calculer les longueurs $AB$ et $AC$.
  \item Calculer le produit scalaire $\vect{AB}\cdot\vect{AC}$.
  \item En déduire une mesure de l'angle $\widehat{BAC}$.
\end{questions}
\reponse{a)~$AB=AC=\sqrt{2}$ ;\; b)~$1$ ;\; c)~$\frac{\pi}{3}$ ($60^{\circ}$).}
\end{methodebox}

\afaire{exercices 10 à 14}

\partiecours{Droites orthogonales, droite orthogonale à un plan}

\begin{minipage}[c]{0.62\linewidth}\raggedright
\begin{definitionbox}[Définition 4 --- Droites orthogonales]
Deux droites de l'espace sont \textbf{orthogonales} si leurs vecteurs directeurs sont
orthogonaux.

Deux droites orthogonales ne sont pas forcément sécantes. Deux droites orthogonales
et sécantes sont dites \textbf{perpendiculaires}.

Dans le cube : $(AB)$ et $(EH)$ sont orthogonales, sans être sécantes ; $(AB)$ et
$(AE)$ sont perpendiculaires.
\end{definitionbox}
\end{minipage}\hfill
\begin{minipage}[c]{0.34\linewidth}\centering
\begin{tikzpicture}
  \cubeespace{2}
  \draw[coulExo,line width=1.2pt] (A)--(B);
  \draw[coulCourbe,line width=1.2pt] (E)--(H);
  \draw[coulCourbe,line width=1.2pt] (A)--(E);
\end{tikzpicture}
\end{minipage}

\begin{minipage}[c]{0.62\linewidth}\raggedright
\begin{definitionbox}[Définition 5 --- Droite orthogonale à un plan]
Une droite est \textbf{orthogonale à un plan} si elle est orthogonale à toutes les
droites de ce plan.
\end{definitionbox}
\end{minipage}\hfill
\begin{minipage}[c]{0.34\linewidth}\centering
\begin{tikzpicture}[font=\footnotesize]
  \fill[gray!15] (0,0)--(3,0)--(3.8,1.1)--(0.8,1.1)--cycle;
  \draw (0,0)--(3,0)--(3.8,1.1)--(0.8,1.1)--cycle;
  \coordinate (O) at (1.9,0.55);
  \draw[coulCourbe,thick] (0.7,0.15)--(1.5,0.95);
  \draw[coulCourbe,thick] (2.3,0.15)--(3.5,0.75);
  \draw[coulCourbe,thick] (2.4,0.97)--(3.6,0.97);
  \draw[coulExo,thick] (O)--(1.9,2.1) node[right]{$(d)$};
  \draw[coulExo,thick,dashed] (O)--(1.9,0);
  \draw[coulExo,thick] (1.9,0)--(1.9,-0.4);
  \path (O) node[croix]{};
  \node at (0.35,0.12) {$\mathcal{P}$};
\end{tikzpicture}

{\footnotesize $(d)$ est orthogonale à toutes les droites de $\mathcal{P}$, même à
celles qu'elle ne coupe pas (en couleur).}
\end{minipage}

\begin{minipage}[c]{0.62\linewidth}\raggedright
\begin{proprietebox}[Propriété 4 --- Caractérisation (admise)]
Une droite est orthogonale à un plan si et seulement si elle est orthogonale à
\textbf{deux droites sécantes} de ce plan.
\end{proprietebox}
\end{minipage}\hfill
\begin{minipage}[c]{0.34\linewidth}\centering
\begin{tikzpicture}[font=\footnotesize]
  \fill[gray!15] (0,0)--(3,0)--(3.8,1.1)--(0.8,1.1)--cycle;
  \draw (0,0)--(3,0)--(3.8,1.1)--(0.8,1.1)--cycle;
  \coordinate (O) at (1.9,0.55);
  \draw[coulCourbe,thick] (0.5,0.3)--(3.3,0.8) node[right]{$(d_{1})$};
  \draw[coulCourbe,thick] (1.1,0.95)--(2.8,0.1) node[below]{$(d_{2})$};
  \draw[coulExo,thick] (O)--(1.9,2.1) node[right]{$(d)$};
  \draw[coulExo,thick,dashed] (O)--(1.9,0);
  \draw[coulExo,thick] (1.9,0)--(1.9,-0.4);
  \draw ($(O)+(0,0.25)$)--++(0.25,0.045)--++(0,-0.25);
  \draw ($(O)+(0,0.25)$)--++(-0.22,0.11)--++(0,-0.25);
  \node at (0.4,0.15) {$\mathcal{P}$};
\end{tikzpicture}

{\footnotesize $(d)\perp(d_{1})$ et $(d)\perp(d_{2})$, sécantes dans $\mathcal{P}$.}
\end{minipage}

\begin{methodebox}[Méthode 4 --- Démontrer que deux droites sont orthogonales]
\etapeexemples
\begin{questions}[label=\alph*)]
  \item \begin{minipage}[t]{0.66\linewidth}\raggedright
        $ABCD$ est un tétraèdre régulier (ses quatre faces sont des triangles
        équilatéraux) d'arête $4$. Démontrer que les droites $(AD)$ et $(BC)$ sont
        orthogonales.
        \end{minipage}\hfill
        \begin{minipage}[t]{0.28\linewidth}\centering\vspace{0pt}
        \begin{tikzpicture}[font=\footnotesize,scale=0.8]
          \coordinate (B) at (0,0); \coordinate (C) at (2.6,0.1);
          \coordinate (D) at (1.6,0.8); \coordinate (A) at (1.3,2.4);
          \draw (A)--(B)--(C)--cycle (A)--(D)--(C);
          \draw[dashed] (B)--(D);
          \path (A) node[above]{$A$} (B) node[left]{$B$} (C) node[right]{$C$}
                (D) node[above right]{$D$};
        \end{tikzpicture}
        \end{minipage}
  \item $ABCDEFGH$ est un cube. Dans le repère orthonormé
        $(A\,;\vect{AB},\vect{AD},\vect{AE})$, démontrer que les droites $(DF)$ et
        $(EB)$ sont orthogonales.
\end{questions}
\begin{correction}
\cas{a)}
\explique{On montre que $\vect{AD}\cdot\vect{BC}=0$. On décompose $\vect{BC}$ en passant
par $A$, pour faire apparaître des angles des faces :}{relation de Chasles}
\begin{calculs}
\vect{AD}\cdot\vect{BC} &= \vect{AD}\cdot\left(\vect{BA}+\vect{AC}\right) &\qquad& \rem{Chasles}\\
\vect{AD}\cdot\vect{BC} &= -\vect{AD}\cdot\vect{AB}+\vect{AD}\cdot\vect{AC} && \rem{bilinéarité, $\vect{BA}=-\vect{AB}$}
\end{calculs}
\explique{Le triangle $ABD$ est équilatéral, donc $\widehat{BAD}=\frac{\pi}{3}$ et
$\vect{AD}\cdot\vect{AB}=4\times 4\times\cos\left(\frac{\pi}{3}\right)=8$.}{face du tétraèdre}
\explique{De même, dans le triangle équilatéral $ACD$ : $\vect{AD}\cdot\vect{AC}=8$.}{}
\explique{Donc $\vect{AD}\cdot\vect{BC}=-8+8=0$ : les vecteurs $\vect{AD}$ et $\vect{BC}$
sont orthogonaux, et \res{les droites $(AD)$ et $(BC)$ sont orthogonales}.}{vecteurs
directeurs orthogonaux}
\cas{b)}\par\nopagebreak
\explique{Dans ce repère : $D(0\,;1\,;0)$, $F(1\,;0\,;1)$, $E(0\,;0\,;1)$ et
$B(1\,;0\,;0)$, d'où $\vect{DF}\begin{pmatrix}1\\-1\\1\end{pmatrix}$ et
$\vect{EB}\begin{pmatrix}1\\0\\-1\end{pmatrix}$.}{on lit les coordonnées}
\begin{minipage}[c]{0.62\linewidth}\raggedright
\begin{calculs}
\vect{DF}\cdot\vect{EB} &= 1\times 1+(-1)\times 0+1\times(-1) &\qquad& \rem{repère orthonormé}\\
\vect{DF}\cdot\vect{EB} &= 0 &&
\end{calculs}
\end{minipage}\hfill
\begin{minipage}[c]{0.34\linewidth}\centering
\begin{tikzpicture}[font=\footnotesize]
  \cubeespace{2}
  \draw[coulExo,thick] (D)--(F) node[pos=0.8,left]{$(DF)$};
  \draw[coulCourbe,thick] (E)--(B) node[pos=0.8,below left]{$(EB)$};
\end{tikzpicture}
\end{minipage}
\explique{Les vecteurs directeurs $\vect{DF}$ et $\vect{EB}$ sont orthogonaux, donc
\res{les droites $(DF)$ et $(EB)$ sont orthogonales}.}{}
\end{correction}

\etape{À vous de jouer}
\begin{questions}[label=\alph*)]
  \item $ABCD$ est un tétraèdre régulier d'arête $6$. Démontrer que les droites $(AB)$
        et $(CD)$ sont orthogonales.
  \item $ABCDEFGH$ est un cube. Dans le repère orthonormé
        $(A\,;\vect{AB},\vect{AD},\vect{AE})$, démontrer que les droites $(HB)$ et
        $(AC)$ sont orthogonales.
\end{questions}
\reponse{a)~$\vect{AB}\cdot\vect{CD}=-18+18=0$ ;\; b)~$\vect{HB}\cdot\vect{AC}=1-1+0=0$.}
\end{methodebox}

\afaire{exercices 15 à 18}

\partiecours{Vecteur normal à un plan}

\begin{minipage}[c]{0.62\linewidth}\raggedright
\begin{definitionbox}[Définition 6 --- Vecteur normal]
Un vecteur non nul $\vect{n}$ est \textbf{normal} à un plan $\mathcal{P}$ s'il dirige
une droite orthogonale à $\mathcal{P}$.
\end{definitionbox}
\end{minipage}\hfill
\begin{minipage}[c]{0.34\linewidth}\centering
\begin{tikzpicture}[font=\footnotesize]
  \fill[gray!15] (0,0)--(3,0)--(3.8,1)--(0.8,1)--cycle;
  \draw (0,0)--(3,0)--(3.8,1)--(0.8,1)--cycle;
  \coordinate (A) at (1.9,0.5);
  \draw[-{Stealth},coulExo,thick] (A)--(1.9,2) node[midway,right]{$\vect{n}$};
  \draw[-{Stealth},coulCourbe,thick] (A)--(2.9,0.5) node[midway,below]{$\vect{u}$};
  \draw[-{Stealth},coulCourbe,thick] (A)--(1.5,0.9) node[midway,left]{$\vect{v}$};
  \node at (0.5,0.2) {$\mathcal{P}$};
\end{tikzpicture}
\end{minipage}

\begin{proprietebox}[Propriété 5 --- Vecteur normal et orthogonalité]
\begin{itemize}
  \item Un vecteur non nul $\vect{n}$ est normal à un plan $\mathcal{P}$ si et
        seulement s'il est orthogonal à \textbf{deux vecteurs non colinéaires}
        $\vect{u}$ et $\vect{v}$ de $\mathcal{P}$ (propriété 4).
  \item Une droite est orthogonale à un plan si et seulement si un vecteur directeur
        de la droite est colinéaire à un vecteur normal du plan.
\end{itemize}
\end{proprietebox}

\begin{methodebox}[Méthode 5 --- Montrer qu'un vecteur est normal à un plan, ou en déterminer un]
\etapeexemples
Dans un repère orthonormé de l'espace :
\begin{questions}[label=\alph*)]
  \item on considère les points $A(1\,;2\,;0)$, $B(3\,;1\,;1)$ et $C(0\,;4\,;1)$.
        Montrer que $A$, $B$ et $C$ définissent un plan, puis que le vecteur
        $\vect{n}\begin{pmatrix}1\\1\\-1\end{pmatrix}$ est normal au plan $(ABC)$ ;
  \item on considère les points $D(1\,;0\,;1)$, $E(2\,;1\,;3)$ et $F(3\,;-1\,;1)$, qui
        définissent un plan. Déterminer un vecteur normal au plan $(DEF)$.
\end{questions}
\begin{correction}
\cas{a)}
\explique{$\vect{AB}\begin{pmatrix}2\\-1\\1\end{pmatrix}$ et $\vect{AC}\begin{pmatrix}-1\\2\\1\end{pmatrix}$ ont
des coordonnées non proportionnelles ($\frac{2}{-1}\neq\frac{1}{1}$) : ils ne sont pas
colinéaires.}{}
\explique{Les points $A$, $B$ et $C$ ne sont pas alignés, donc ils définissent un
plan.}{}
\explique{Le repère est orthonormé :}{}
\begin{calculs}
\vect{n}\cdot\vect{AB} &= 1\times 2+1\times(-1)+(-1)\times 1 &\qquad& \\
\vect{n}\cdot\vect{AB} &= 0 &&
\end{calculs}
\begin{calculs}
\vect{n}\cdot\vect{AC} &= 1\times(-1)+1\times 2+(-1)\times 1 &\qquad& \\
\vect{n}\cdot\vect{AC} &= 0 &&
\end{calculs}
\explique{$\vect{n}$ est orthogonal à deux vecteurs non colinéaires du plan $(ABC)$,
donc \res{$\vect{n}$ est normal au plan $(ABC)$}.}{propriété 5}
\cas{b)}
\explique{On a $\vect{DE}\begin{pmatrix}1\\1\\2\end{pmatrix}$ et $\vect{DF}\begin{pmatrix}2\\-1\\0\end{pmatrix}$, non
colinéaires. On cherche un vecteur non nul $\vect{n}\begin{pmatrix}a\\b\\c\end{pmatrix}$ orthogonal à $\vect{DE}$ et à
$\vect{DF}$ :}{deux équations, trois inconnues}
\begin{calculs}
& \accol{\vect{n}\cdot\vect{DE} &=& 0\\ \vect{n}\cdot\vect{DF} &=& 0} &\qquad& \\
\iff & \accol{a+b+2c &=& 0\\ 2a-b &=& 0} && \rem{repère orthonormé}\\
\iff & \accol{a+2a+2c &=& 0\\ b &=& 2a} && \rem{on remplace $b$ par $2a$}\\
\iff & \accol{c &=& -\frac{3}{2}a\\ b &=& 2a} &&
\end{calculs}
\explique{On choisit $a=2$ (pour éviter les fractions) : $b=4$ et $c=-3$.}{un choix
parmi une infinité}
\explique{Le vecteur \res{$\vect{n}\begin{pmatrix}2\\4\\-3\end{pmatrix}$} est normal au plan $(DEF)$.}{tout
vecteur colinéaire non nul convient aussi}
\end{correction}

\etape{À vous de jouer}
Dans un repère orthonormé de l'espace :
\begin{questions}[label=\alph*)]
  \item on considère les points $A(2\,;0\,;1)$, $B(3\,;2\,;1)$ et $C(1\,;1\,;3)$.
        Montrer que $A$, $B$ et $C$ définissent un plan, puis que le vecteur
        $\vect{n}\begin{pmatrix}4\\-2\\3\end{pmatrix}$ est normal au plan $(ABC)$ ;
  \item on considère les points $D(0\,;1\,;2)$, $E(1\,;0\,;3)$ et $F(0\,;3\,;3)$, qui
        définissent un plan. Déterminer un vecteur normal au plan $(DEF)$.
\end{questions}
\reponse{a)~$\vect{n}\cdot\vect{AB}=\vect{n}\cdot\vect{AC}=0$ ;\; b)~par exemple
$\vect{n}\begin{pmatrix}3\\1\\-2\end{pmatrix}$.}
\end{methodebox}

\afaire{exercices 19 à 22}

\partiecours{Équation cartésienne d'un plan}

\begin{proprietebox}[Théorème 1 --- Équation cartésienne d'un plan]
L'espace est muni d'un repère orthonormé.
\begin{itemize}
  \item Le plan $\mathcal{P}$ passant par un point $A$ et de vecteur normal
        $\vect{n}\begin{pmatrix}a\\b\\c\end{pmatrix}$ est l'ensemble des points $M$ tels que
        $\vect{AM}\cdot\vect{n}=0$. Il a une équation de la forme
        \[
        ax+by+cz+d=0,\quad\text{où }d\text{ est un nombre réel.}
        \]
  \item Réciproquement, si $a$, $b$, $c$ ne sont pas tous nuls, l'ensemble des points
        $M(x\,;y\,;z)$ tels que $ax+by+cz+d=0$ est un plan, de vecteur normal
        $\vect{n}\begin{pmatrix}a\\b\\c\end{pmatrix}$.
\end{itemize}
\end{proprietebox}

\textbf{\color{coulMeth}Démonstration} (exigible) du premier point.
Soit $A(x_{A}\,;y_{A}\,;z_{A})$ et $M(x\,;y\,;z)$ un point de l'espace. Le point $M$
est dans $\mathcal{P}$ si et seulement si $\vect{AM}$ est un vecteur de $\mathcal{P}$,
c'est-à-dire orthogonal au vecteur normal $\vect{n}$ : $\vect{AM}\cdot\vect{n}=0$.

\smallskip
On calcule ce produit scalaire avec $\vect{AM}\begin{pmatrix}x-x_{A}\\y-y_{A}\\z-z_{A}\end{pmatrix}$ :
\begin{calculs}
\vect{AM}\cdot\vect{n} &= a(x-x_{A})+b(y-y_{A})+c(z-z_{A}) &\qquad& \rem{repère orthonormé}\\
\vect{AM}\cdot\vect{n} &= ax+by+cz-(ax_{A}+by_{A}+cz_{A}) && \rem{on développe}
\end{calculs}
On pose $d=-(ax_{A}+by_{A}+cz_{A})$, qui ne dépend que de $A$ et de $\vect{n}$. Donc
$M\in\mathcal{P}$ si et seulement si $ax+by+cz+d=0$.

\begin{methodebox}[Méthode 6 --- Déterminer une équation cartésienne d'un plan]
\etapeexemples
Dans un repère orthonormé de l'espace, déterminer une équation cartésienne :
\begin{questions}[label=\alph*)]
  \item du plan $\mathcal{P}$ passant par $A(2\,;-1\,;3)$ et de vecteur normal
        $\vect{n}\begin{pmatrix}1\\-2\\4\end{pmatrix}$ ;
  \item du plan $(ABC)$, avec $A(1\,;2\,;0)$, $B(3\,;1\,;1)$, $C(0\,;4\,;1)$ (méthode 5 :
        $\vect{n}\begin{pmatrix}1\\1\\-1\end{pmatrix}$ est normal à $(ABC)$) ;
  \item du plan médiateur $\mathcal{Q}$ du segment $[EF]$, avec $E(1\,;3\,;-2)$ et
        $F(3\,;-1\,;4)$ : c'est le plan qui passe par le milieu de $[EF]$ et qui est
        orthogonal à $(EF)$.
\end{questions}
\begin{correction}
\cas{a)}
\explique{$\vect{n}\begin{pmatrix}1\\-2\\4\end{pmatrix}$ est normal à $\mathcal{P}$, donc $\mathcal{P}$ a une équation
de la forme $x-2y+4z+d=0$, où $d$ est un nombre réel.}{théorème 1}
\explique{$A\in\mathcal{P}$, donc ses coordonnées vérifient l'équation :}{on trouve $d$}
\begin{calculs}
2-2\times(-1)+4\times 3+d &= 0 &\qquad& \\
\iff 16+d &= 0 &&\\
\iff d &= -16 &&
\end{calculs}
\explique{Une équation cartésienne de $\mathcal{P}$ est \res{$x-2y+4z-16=0$}.}{}
\cas{b)}
\explique{$(ABC)$ a une équation de la forme $x+y-z+d=0$. $A\in(ABC)$, donc
$1+2-0+d=0$, d'où $d=-3$.}{même démarche}
\explique{Une équation cartésienne de $(ABC)$ est \res{$x+y-z-3=0$}.}{contrôle avec $B$ :
$3+1-1-3=0$}
\cas{c)}
\explique{Le milieu de $[EF]$ est $I\left(\frac{1+3}{2}\,;\frac{3-1}{2}\,;\frac{-2+4}{2}\right)$,
soit $I(2\,;1\,;1)$.}{un point du plan}
\explique{$\vect{EF}\begin{pmatrix}2\\-4\\6\end{pmatrix}$ dirige une droite orthogonale à $\mathcal{Q}$ : il est normal à
$\mathcal{Q}$, comme $\frac{1}{2}\vect{EF}\begin{pmatrix}1\\-2\\3\end{pmatrix}$.}{un vecteur normal, simplifié}
\explique{$\mathcal{Q}$ a une équation de la forme $x-2y+3z+d=0$ ; $I\in\mathcal{Q}$,
donc $2-2+3+d=0$, d'où $d=-3$.}{}
\explique{Une équation cartésienne de $\mathcal{Q}$ est \res{$x-2y+3z-3=0$}.}{}
\end{correction}

\etape{À vous de jouer}
Dans un repère orthonormé de l'espace, déterminer une équation cartésienne :
\begin{questions}[label=\alph*)]
  \item du plan $\mathcal{P}$ passant par $A(-1\,;2\,;1)$ et de vecteur normal
        $\vect{n}\begin{pmatrix}3\\-1\\2\end{pmatrix}$ ;
  \item du plan $(ABC)$, avec $A(2\,;0\,;1)$, $B(3\,;2\,;1)$, $C(1\,;1\,;3)$ (« À vous de
        jouer » de la méthode 5 : $\vect{n}\begin{pmatrix}4\\-2\\3\end{pmatrix}$ est normal à $(ABC)$) ;
  \item du plan médiateur $\mathcal{Q}$ du segment $[EF]$, avec $E(0\,;2\,;1)$ et
        $F(4\,;-2\,;3)$.
\end{questions}
\reponse{a)~$3x-y+2z+3=0$ ;\; b)~$4x-2y+3z-11=0$ ;\; c)~$2x-2y+z-6=0$.}
\end{methodebox}

\begin{proprietebox}[Propriété 6 --- Positions relatives d'une droite et d'un plan]
Soit $(d)$ une droite passant par $A$, de vecteur directeur $\vect{u}$, et $\mathcal{P}$
un plan de vecteur normal $\vect{n}$. Trois cas sont possibles.

\smallskip
\begin{minipage}[t]{0.32\linewidth}\centering
\begin{tikzpicture}[font=\footnotesize]
  \fill[gray!15,shift={(0,0)}] (0,0)--(2.6,0)--(3.3,0.9)--(0.7,0.9)--cycle;
  \draw[shift={(0,0)}] (0,0)--(2.6,0)--(3.3,0.9)--(0.7,0.9)--cycle;
  \node[shift={(0,0)}] at (2.2,0.17) {$\mathcal{P}$};
  \draw[-{Stealth},thick] (0.9,0.3)--(0.9,1.0) node[midway,left]{$\vect{n}$};
  \draw[coulExo,thick] (1.9,0.45)--(2.4,1.9) node[right]{$(d)$};
  \draw[coulExo,thick,dashed] (1.9,0.45)--(1.75,0);
  \draw[coulExo,thick] (1.75,0)--(1.65,-0.3);
  \path (1.9,0.45) node[croix]{};
\end{tikzpicture}

\textbf{Sécants} en un point\\
$\vect{u}\cdot\vect{n}\neq 0$
\end{minipage}\hfill
\begin{minipage}[t]{0.32\linewidth}\centering
\begin{tikzpicture}[font=\footnotesize]
  \fill[gray!15,shift={(0,0)}] (0,0)--(2.6,0)--(3.3,0.9)--(0.7,0.9)--cycle;
  \draw[shift={(0,0)}] (0,0)--(2.6,0)--(3.3,0.9)--(0.7,0.9)--cycle;
  \node[shift={(0,0)}] at (2.2,0.17) {$\mathcal{P}$};
  \draw[-{Stealth},thick] (0.9,0.3)--(0.9,1.0) node[midway,left]{$\vect{n}$};
  \draw[coulExo,thick] (0.6,1.5)--(3.2,1.7) node[above]{$(d)$};
  \path (1.4,1.56) node[croix]{} node[above]{$A$};
  \draw (1.5,-0.3) node{};
\end{tikzpicture}

\textbf{Strictement parallèles}\\
$\vect{u}\cdot\vect{n}=0$ et $A\notin\mathcal{P}$
\end{minipage}\hfill
\begin{minipage}[t]{0.32\linewidth}\centering
\begin{tikzpicture}[font=\footnotesize]
  \fill[gray!15,shift={(0,0)}] (0,0)--(2.6,0)--(3.3,0.9)--(0.7,0.9)--cycle;
  \draw[shift={(0,0)}] (0,0)--(2.6,0)--(3.3,0.9)--(0.7,0.9)--cycle;
  \node[shift={(0,0)}] at (2.2,0.17) {$\mathcal{P}$};
  \draw[-{Stealth},thick] (0.9,0.3)--(0.9,1.0) node[midway,left]{$\vect{n}$};
  \draw[coulExo,thick] (0.4,0.25)--(3.0,0.65) node[above left]{$(d)$};
  \path (1.6,0.43) node[croix]{} node[below]{$A$};
  \draw (1.5,-0.3) node{};
\end{tikzpicture}

\textbf{Droite contenue dans le plan}\\
$\vect{u}\cdot\vect{n}=0$ et $A\in\mathcal{P}$
\end{minipage}

\smallskip
En effet, $\vect{u}\cdot\vect{n}=0$ signifie que $\vect{u}$ est un vecteur de la
direction de $\mathcal{P}$ : la droite est alors parallèle au plan.
\end{proprietebox}

\begin{proprietebox}[Propriété 7 --- Positions relatives de deux plans]
Soit $\mathcal{P}$ et $\mathcal{P}'$ deux plans de vecteurs normaux $\vect{n}$ et
$\vect{n'}$, et $A$ un point de $\mathcal{P}$. Trois cas sont possibles.

\smallskip
\begin{minipage}[t]{0.32\linewidth}\centering
\begin{tikzpicture}[font=\footnotesize]
  \fill[gray!15,shift={(0,0)}] (0,0)--(2.6,0)--(3.3,0.9)--(0.7,0.9)--cycle;
  \draw[shift={(0,0)}] (0,0)--(2.6,0)--(3.3,0.9)--(0.7,0.9)--cycle;
  \node[shift={(0,0)}] at (2.2,0.17) {$\mathcal{P}$};
  \fill[coulExo,opacity=0.15] (0.35,-0.5)--(2.95,-0.5)--(2.95,1.4)--(0.35,1.4)--cycle;
  \draw[coulExo] (0.35,-0.5)--(2.95,-0.5)--(2.95,1.4)--(0.35,1.4)--cycle;
  \draw[thick] (0.35,0.45)--(2.95,0.45);
  \node[coulExo] at (2.6,1.15) {$\mathcal{P}'$};
\end{tikzpicture}

\textbf{Sécants} suivant une droite\\
$\vect{n}$ et $\vect{n'}$ non colinéaires
\end{minipage}\hfill
\begin{minipage}[t]{0.32\linewidth}\centering
\begin{tikzpicture}[font=\footnotesize]
  \fill[gray!15,shift={(0,0)}] (0,0)--(2.6,0)--(3.3,0.9)--(0.7,0.9)--cycle;
  \draw[shift={(0,0)}] (0,0)--(2.6,0)--(3.3,0.9)--(0.7,0.9)--cycle;
  \node[shift={(0,0)}] at (2.2,0.17) {$\mathcal{P}$};
  \fill[gray!15,shift={(0,1.1)}] (0,0)--(2.6,0)--(3.3,0.9)--(0.7,0.9)--cycle;
  \draw[shift={(0,1.1)}] (0,0)--(2.6,0)--(3.3,0.9)--(0.7,0.9)--cycle;
  \node[shift={(0,1.1)}] at (2.2,0.17) {$\mathcal{P}'$};
  \path (1.2,0.5) node[croix]{} node[below]{$A$};
\end{tikzpicture}

\textbf{Strictement parallèles}\\
$\vect{n}$ et $\vect{n'}$ colinéaires et $A\notin\mathcal{P}'$
\end{minipage}\hfill
\begin{minipage}[t]{0.32\linewidth}\centering
\begin{tikzpicture}[font=\footnotesize]
  \fill[gray!15,shift={(0,0)}] (0,0)--(2.6,0)--(3.3,0.9)--(0.7,0.9)--cycle;
  \draw[shift={(0,0)}] (0,0)--(2.6,0)--(3.3,0.9)--(0.7,0.9)--cycle;
  \node[shift={(0,0)}] at (2.2,0.17) {$\mathcal{P}=\mathcal{P}'$};
  \path (1.2,0.5) node[croix]{} node[below]{$A$};
  \draw (1.5,1.6) node{};
\end{tikzpicture}

\textbf{Confondus}\\
$\vect{n}$ et $\vect{n'}$ colinéaires et $A\in\mathcal{P}'$
\end{minipage}

\smallskip
Deux plans sont dits \textbf{perpendiculaires} si $\vect{n}\cdot\vect{n'}=0$ ; ils
sont alors sécants (premier cas).
\end{proprietebox}

\begin{methodebox}[Méthode 7 --- Étudier la position relative d'une droite et d'un plan]
\etapeexemples
Dans un repère orthonormé, on considère le plan $\mathcal{P}\colon 2x-y+z-3=0$ et les
droites
\[
(d)\colon\left\{\begin{array}{l} x=1+t\\ y=-1+2t\\ z=2-t\end{array}\right.
\qquad
(d')\colon\left\{\begin{array}{l} x=2t\\ y=1+t\\ z=-3t\end{array}\right.
\qquad (t\in\R).
\]
\begin{questions}[label=\alph*)]
  \item Montrer que $(d)$ et $\mathcal{P}$ sont sécants et déterminer leur point
        d'intersection.
  \item Étudier la position relative de $(d')$ et de $\mathcal{P}$.
\end{questions}
\begin{correction}
\cas{a)}
\explique{$\vect{n}\begin{pmatrix}2\\-1\\1\end{pmatrix}$ est normal à $\mathcal{P}$ et $\vect{u}\begin{pmatrix}1\\2\\-1\end{pmatrix}$ dirige
$(d)$.}{on lit les coefficients}
\explique{$\vect{u}\cdot\vect{n}=1\times 2+2\times(-1)+(-1)\times 1=-1$, donc
$\vect{u}\cdot\vect{n}\neq 0$ : $(d)$ et $\mathcal{P}$ sont sécants.}{propriété 6}
\explique{Le point d'intersection est sur $(d)$ : il existe $t\in\R$ tel que ses
coordonnées sont $(1+t\,;-1+2t\,;2-t)$. Il est sur $\mathcal{P}$, donc elles vérifient
l'équation de $\mathcal{P}$ :}{}
\begin{calculs}
2(1+t)-(-1+2t)+(2-t)-3 &= 0 &\qquad& \rem{on remplace $x$, $y$, $z$}\\
\iff 2+2t+1-2t+2-t-3 &= 0 &&\\
\iff 2-t &= 0 &&\\
\iff t &= 2 &&
\end{calculs}
\explique{Pour $t=2$ : $x=1+2=3$, $y=-1+4=3$ et $z=2-2=0$.}{contrôle :
$6-3+0-3=0$}
\explique{$(d)$ et $\mathcal{P}$ sont sécants en \res{$K(3\,;3\,;0)$}.}{}
\cas{b)}
\explique{$\vect{u'}\begin{pmatrix}2\\1\\-3\end{pmatrix}$ dirige $(d')$ et
$\vect{u'}\cdot\vect{n}=4-1-3=0$ : $(d')$ est parallèle à $\mathcal{P}$.}{propriété 6}
\explique{Le point $(0\,;1\,;0)$ de $(d')$ (pour $t=0$) donne $2\times 0-1+0-3=-4\neq 0$ :
il n'est pas dans $\mathcal{P}$.}{on teste un point}
\explique{\res{$(d')$ est strictement parallèle à $\mathcal{P}$}.}{}
\end{correction}

\etape{À vous de jouer}
Dans un repère orthonormé, on considère le plan $\mathcal{P}\colon x+2y-z+1=0$ et les
droites
\[
(d)\colon\left\{\begin{array}{l} x=2-t\\ y=t\\ z=1+3t\end{array}\right.
\qquad
(d')\colon\left\{\begin{array}{l} x=1+t\\ y=t\\ z=2+3t\end{array}\right.
\qquad (t\in\R).
\]
\begin{questions}[label=\alph*)]
  \item Montrer que $(d)$ et $\mathcal{P}$ sont sécants et déterminer leur point
        d'intersection.
  \item Étudier la position relative de $(d')$ et de $\mathcal{P}$.
\end{questions}
\reponse{a)~Sécants en $(1\,;1\,;4)$ ;\; b)~$(d')$ est contenue dans $\mathcal{P}$.}
\end{methodebox}

\begin{methodebox}[Méthode 8 --- Étudier la position relative de deux plans]
\etapeexemples
Dans un repère orthonormé, on considère les plans
\begin{center}
$\mathcal{P}\colon 2x+y-2z+1=0$ \qquad $\mathcal{Q}\colon x-4y-z+3=0$ \qquad
$\mathcal{R}\colon -4x-2y+4z+5=0$
\end{center}
\begin{questions}[label=\alph*)]
  \item Montrer que les plans $\mathcal{P}$ et $\mathcal{Q}$ sont perpendiculaires.
  \item Montrer que les plans $\mathcal{P}$ et $\mathcal{R}$ sont strictement parallèles.
\end{questions}
\begin{correction}
\cas{a)}
\explique{$\vect{n}\begin{pmatrix}2\\1\\-2\end{pmatrix}$ est normal à $\mathcal{P}$ et $\vect{n'}\begin{pmatrix}1\\-4\\-1\end{pmatrix}$ est
normal à $\mathcal{Q}$.}{coefficients de $x$, $y$, $z$}
\begin{calculs}
\vect{n}\cdot\vect{n'} &= 2\times 1+1\times(-4)+(-2)\times(-1) &\qquad& \\
\vect{n}\cdot\vect{n'} &= 0 &&
\end{calculs}
\explique{Les vecteurs normaux sont orthogonaux, donc \res{$\mathcal{P}$ et
$\mathcal{Q}$ sont perpendiculaires}.}{propriété 7}
\cas{b)}
\explique{$\vect{n''}\begin{pmatrix}-4\\-2\\4\end{pmatrix}$ est normal à $\mathcal{R}$ et $\vect{n''}=-2\vect{n}$ : les
vecteurs normaux sont colinéaires, donc $\mathcal{P}$ et $\mathcal{R}$ sont
parallèles.}{propriété 7}
\explique{Le point $(0\,;-1\,;0)$ est dans $\mathcal{P}$ ($0-1-0+1=0$), mais pas dans
$\mathcal{R}$ ($0+2+0+5=7\neq 0$).}{un point de l'un, testé dans l'autre}
\explique{Les plans n'ont pas de point commun : \res{$\mathcal{P}$ et $\mathcal{R}$ sont
strictement parallèles}.}{}
\end{correction}

\etape{À vous de jouer}
Dans un repère orthonormé, on considère les plans
\begin{center}
$\mathcal{P}\colon x-3y+2z-4=0$ \qquad $\mathcal{Q}\colon x+y+z+2=0$ \qquad
$\mathcal{R}\colon 2x-6y+4z-8=0$
\end{center}
\begin{questions}[label=\alph*)]
  \item Montrer que les plans $\mathcal{P}$ et $\mathcal{Q}$ sont perpendiculaires.
  \item Étudier la position relative des plans $\mathcal{P}$ et $\mathcal{R}$.
\end{questions}
\reponse{a)~$\vect{n}\cdot\vect{n'}=1-3+2=0$ ;\; b)~$\mathcal{P}$ et $\mathcal{R}$ sont
confondus.}
\end{methodebox}

\afaire{exercices 23 à 31}

\partiecours{Projeté orthogonal et distances}

\begin{minipage}[c]{0.62\linewidth}\raggedright
\begin{definitionbox}[Définition 7 --- Projeté orthogonal]
Soit $M$ un point de l'espace.
\begin{itemize}
  \item Le \textbf{projeté orthogonal} de $M$ sur une droite $(d)$ est le point $H$
        de $(d)$ tel que $(MH)$ est perpendiculaire à $(d)$ ($H=M$ si $M\in(d)$).
  \item Le \textbf{projeté orthogonal} de $M$ sur un plan $\mathcal{P}$ est le point
        $H$ de $\mathcal{P}$ tel que $(MH)$ est orthogonale à $\mathcal{P}$
        ($H=M$ si $M\in\mathcal{P}$).
\end{itemize}
\end{definitionbox}
\end{minipage}\hfill
\begin{minipage}[c]{0.34\linewidth}\centering
\begin{tikzpicture}[font=\footnotesize]
  \fill[gray!15] (0,0)--(3,0)--(3.8,1)--(0.8,1)--cycle;
  \draw (0,0)--(3,0)--(3.8,1)--(0.8,1)--cycle;
  \coordinate (H) at (1.6,0.5); \coordinate (M) at (1.6,2.1);
  \coordinate (N) at (2.9,0.6);
  \draw[thick] (M)--(H);
  \draw[dashed] (M)--(N) (H)--(N);
  \draw (H)++(0,0.2)--++(0.2,0)--++(0,-0.2);
  \path (M) node[croix]{} node[above]{$M$} (H) node[croix]{} node[below left]{$H$}
        (N) node[croix]{} node[below right]{$N$};
  \node at (0.5,0.2) {$\mathcal{P}$};
\end{tikzpicture}

\medskip
\begin{tikzpicture}[font=\footnotesize]
  \coordinate (H) at (1.6,0.4); \coordinate (M) at (1.9,1.6);
  \draw[coulExo,thick] (0,0)--(3.6,0.9) node[above]{$(d)$};
  \draw[thick] (M)--(H);
  \draw ($(H)!0.2!(M)$)--($(H)!0.2!(M)+(0.194,0.0485)$)--($(H)+(0.194,0.0485)$);
  \path (M) node[croix]{} node[above]{$M$} (H) node[croix]{} node[below right]{$H$};
\end{tikzpicture}
\end{minipage}

\begin{proprietebox}[Propriété 8 --- Le point le plus proche]
Le projeté orthogonal $H$ d'un point $M$ sur un plan $\mathcal{P}$ (ou sur une droite)
est le point de $\mathcal{P}$ (de la droite) le plus proche de $M$. La distance $MH$
est appelée \textbf{distance du point $M$ au plan $\mathcal{P}$} (à la droite).
\end{proprietebox}

\textbf{\color{coulMeth}Démonstration} (exigible), pour un plan.
Soit $N$ un point de $\mathcal{P}$, distinct de $H$. Le vecteur $\vect{HN}$ est un
vecteur du plan $\mathcal{P}$ et $(MH)$ est orthogonale à $\mathcal{P}$, donc
$\vect{MH}\cdot\vect{HN}=0$.
\begin{calculs}
MN^{2} &= \|\vect{MH}+\vect{HN}\|^{2} &\qquad& \rem{Chasles}\\
MN^{2} &= MH^{2}+2\,\vect{MH}\cdot\vect{HN}+HN^{2} && \rem{propriétés 1 et 2}\\
MN^{2} &= MH^{2}+HN^{2} && \rem{$\vect{MH}\cdot\vect{HN}=0$}
\end{calculs}
Or $HN^{2}>0$, car $N\neq H$, donc $MN^{2}>MH^{2}$, puis $MN>MH$ (des longueurs
positives). Tout point de $\mathcal{P}$ autre que $H$ est donc plus loin de $M$ que $H$.

\begin{methodebox}[Méthode 9 --- Déterminer le projeté orthogonal d'un point sur un plan et la distance au plan]
\etapeexemples
Dans un repère orthonormé, on considère le plan $\mathcal{P}\colon x-2y+2z-1=0$ et le
point $M(2\,;-1\,;3)$.
\begin{questions}[label=\alph*)]
  \item Déterminer une représentation paramétrique de la droite $\Delta$ passant par
        $M$ et orthogonale à $\mathcal{P}$.
  \item En déduire les coordonnées du projeté orthogonal $H$ de $M$ sur $\mathcal{P}$.
  \item Calculer la distance du point $M$ au plan $\mathcal{P}$.
\end{questions}
\begin{correction}
\cas{a)}
\explique{$\Delta$ est orthogonale à $\mathcal{P}$, donc elle est dirigée par le
vecteur normal $\vect{n}\begin{pmatrix}1\\-2\\2\end{pmatrix}$ de $\mathcal{P}$ ; elle passe par $M$.}{propriété 5}
\explique{Une représentation paramétrique de $\Delta$ est
\res{$\left\{\begin{array}{l} x=2+t\\ y=-1-2t\\ z=3+2t\end{array}\right.$} $(t\in\R)$.}{ch. 4}
\cas{b)}
\explique{$H$ est le point d'intersection de $\Delta$ et de $\mathcal{P}$. Il est sur
$\Delta$ : il existe $t\in\R$ tel que ses coordonnées sont
$(2+t\,;-1-2t\,;3+2t)$. Il est sur $\mathcal{P}$ :}{définition 7}
\begin{calculs}
(2+t)-2(-1-2t)+2(3+2t)-1 &= 0 &\qquad& \rem{on remplace $x$, $y$, $z$}\\
\iff 2+t+2+4t+6+4t-1 &= 0 &&\\
\iff 9t+9 &= 0 &&\\
\iff t &= -1 &&
\end{calculs}
\explique{Pour $t=-1$ : $x=2-1=1$, $y=-1+2=1$ et $z=3-2=1$, donc
\res{$H(1\,;1\,;1)$}.}{contrôle : $1-2+2-1=0$}
\cas{c)}
\explique{La distance de $M$ à $\mathcal{P}$ est $MH$, avec $\vect{MH}\begin{pmatrix}-1\\2\\-2\end{pmatrix}$ :}{propriété 8}
\begin{calculs}
MH &= \sqrt{(-1)^{2}+2^{2}+(-2)^{2}} &\qquad& \rem{repère orthonormé}\\
MH &= \sqrt{9} &&\\
MH &= \res{3} &&
\end{calculs}
\end{correction}

\etape{À vous de jouer}
Dans un repère orthonormé, on considère le plan $\mathcal{P}\colon 2x+y-2z+4=0$ et le
point $M(5\,;2\,;-1)$.
\begin{questions}[label=\alph*)]
  \item Déterminer une représentation paramétrique de la droite $\Delta$ passant par
        $M$ et orthogonale à $\mathcal{P}$.
  \item En déduire les coordonnées du projeté orthogonal $H$ de $M$ sur $\mathcal{P}$.
  \item Calculer la distance du point $M$ au plan $\mathcal{P}$.
\end{questions}
\reponse{a)~$x=5+2t$, $y=2+t$, $z=-1-2t$ ($t\in\R$) ;\; b)~$H(1\,;0\,;3)$ ;\;
c)~$6$.}
\end{methodebox}

\begin{methodebox}[Méthode 10 --- Déterminer le projeté orthogonal d'un point sur une droite et la distance à la droite]
\etapeexemples
Dans un repère orthonormé, on considère la droite $(d)$ passant par $A(1\,;0\,;2)$ et
de vecteur directeur $\vect{u}\begin{pmatrix}2\\1\\-1\end{pmatrix}$, et le point $M(5\,;-1\,;-3)$.
\begin{questions}[label=\alph*)]
  \item Déterminer les coordonnées du projeté orthogonal $H$ de $M$ sur $(d)$.
  \item Calculer la distance du point $M$ à la droite $(d)$.
\end{questions}
\begin{correction}
\cas{a)}
\explique{Une représentation paramétrique de $(d)$ est
$\left\{\begin{array}{l} x=1+2t\\ y=t\\ z=2-t\end{array}\right.$ $(t\in\R)$.}{ch. 4}
\explique{$H\in(d)$ : il existe $t\in\R$ tel que $H(1+2t\,;t\,;2-t)$, d'où
$\vect{MH}\begin{pmatrix}-4+2t\\1+t\\5-t\end{pmatrix}$.}{coordonnées de $H$ moins celles de $M$}
\explique{$(MH)$ est perpendiculaire à $(d)$, donc $\vect{MH}\cdot\vect{u}=0$ :}{définition 7}
\begin{calculs}
2(-4+2t)+1\times(1+t)+(-1)(5-t) &= 0 &\qquad& \rem{repère orthonormé}\\
\iff -8+4t+1+t-5+t &= 0 &&\\
\iff 6t-12 &= 0 &&\\
\iff t &= 2 &&
\end{calculs}
\explique{Pour $t=2$ : $x=1+4=5$, $y=2$ et $z=2-2=0$, donc \res{$H(5\,;2\,;0)$}.}{}
\cas{b)}
\explique{$\vect{MH}\begin{pmatrix}0\\3\\3\end{pmatrix}$ (contrôle : $\vect{MH}\cdot\vect{u}=0+3-3=0$).}{propriété 8}
\begin{calculs}
MH &= \sqrt{0^{2}+3^{2}+3^{2}} &\qquad& \\
MH &= \sqrt{18} &&\\
MH &= \res{3\sqrt{2}} &&
\end{calculs}
\end{correction}

\etape{À vous de jouer}
Dans un repère orthonormé, on considère la droite $(d)$ passant par $A(0\,;1\,;-1)$ et
de vecteur directeur $\vect{u}\begin{pmatrix}1\\-2\\2\end{pmatrix}$, et le point $M(0\,;-5\,;2)$.
\begin{questions}[label=\alph*)]
  \item Déterminer les coordonnées du projeté orthogonal $H$ de $M$ sur $(d)$.
  \item Calculer la distance du point $M$ à la droite $(d)$.
\end{questions}
\reponse{a)~$H(2\,;-3\,;3)$ ;\; b)~$3$.}
\end{methodebox}

\afaire{exercices 32 à 36, puis les exercices type bac 37 à 39}

\end{document}
